\documentclass[11pt,a4paper]{article}

% ---------------------------------------------------------
% PREAMBLE (stable, warning-free, Overleaf-safe)
% ---------------------------------------------------------
\usepackage[utf8]{inputenc}
\usepackage[T1]{fontenc}
\usepackage{lmodern}
\usepackage{amsmath, amssymb, amsfonts}
\usepackage{bm}
\usepackage{geometry}
\usepackage{graphicx}
\usepackage{hyperref}
\usepackage{cite}
\usepackage{authblk}
\usepackage{physics}
\usepackage{mathtools}

\geometry{margin=2.4cm}

\title{\textbf{Companion LXXIX: Score-Matching and Diffusion Models for Persistence Diagrams in the Relaxation-Driven Cyclic Cosmology}}
\author{Michael Lehmann \\ \texttt{mi.lehmann@gmx.de}}
\date{April 2026}

\begin{document}
\maketitle

\begin{abstract}
Score-based diffusion models provide a powerful generative framework for learning 
complex probability distributions through stochastic differential equations. 
Applied to persistence diagrams, these models learn the gradient of the log-density 
(score function), enabling sampling, denoising, and likelihood-free inference in 
topological data analysis. Building on the Schrödinger bridge dynamics introduced 
in Companion LXXVIII, this work develops score-matching and diffusion models for 
persistence diagrams in the Relaxation-Driven Cyclic Cosmology (RDCC). The 
relaxation mechanism modifies the scale-dependent evolution of the gravitational 
potential, inducing characteristic changes in score fields, diffusion trajectories, 
and generative flows. The resulting framework provides a modern, flexible, and 
probabilistically grounded approach to modeling RDCC signatures in weak-lensing 
topology.
\end{abstract}
% ---------------------------------------------------------
\section{Introduction}

Persistence diagrams encode the birth and death of topological features in weak-
lensing convergence fields. Previous Companions introduced Wasserstein geometry, 
entropic optimal transport, and Schrödinger bridges as tools for analyzing the 
Relaxation-Driven Cyclic Cosmology (RDCC).

Score-based diffusion models extend this framework by learning the gradient of the 
log-density of persistence diagrams. This score function governs the reverse-time 
diffusion process that generates new diagrams from noise. In cosmology, this 
provides a powerful method for modeling the distribution of topological features, 
denoising weak-lensing maps, and performing likelihood-free inference.

In RDCC, the relaxation mechanism modifies the underlying distribution of 
persistence diagrams, producing characteristic signatures in the score field and 
the associated diffusion dynamics.
% ---------------------------------------------------------
\section{Score Matching for Persistence Diagrams}

Let $p(D)$ be the probability density of persistence diagrams. The score function 
is

\begin{equation}
    s(D)
    = \nabla_{D} \log p(D).
\end{equation}

Score matching estimates $s(D)$ by minimizing

\begin{equation}
    \mathcal{L}_{\mathrm{SM}}
    = \mathbb{E}_{p}
      \left[
        \| s_{\theta}(D) \|^{2}
        + 2 \, \nabla \cdot s_{\theta}(D)
      \right],
\end{equation}

where $s_{\theta}$ is a neural approximation of the score.

For persistence diagrams, the gradient acts on birth-death coordinates.
% ---------------------------------------------------------
\section{Diffusion Models for Persistence Diagrams}

Forward diffusion adds noise:

\begin{equation}
    dD_{t}
    = \sqrt{\beta(t)} \, dB_{t},
\end{equation}

where $B_{t}$ is Brownian motion on diagram space.

The reverse-time SDE is

\begin{equation}
    dD_{t}
    = \left[
        - \beta(t) s(D_{t}, t)
      \right] dt
      + \sqrt{\beta(t)} \, d\bar{B}_{t},
\end{equation}

where $s(D_{t}, t)$ is the time-dependent score.

Sampling proceeds by integrating this reverse SDE from noise to data.
% ---------------------------------------------------------
\section{RDCC Modification of Score Fields}

In RDCC, the convergence evolves as

\begin{equation}
    \kappa_{\mathrm{RDCC}}(\ell)
    = \kappa(\ell)
      \left[
        1 - \chi_{\kappa}
        \left(\frac{\ell}{\ell_{\mathrm{rel}}}\right)^{\alpha}
      \right].
\end{equation}

This modifies the log-density of persistence diagrams:

\begin{equation}
    \log p_{\mathrm{RDCC}}(D)
    = \log p(D)
      - \chi_{p}
        \left(\frac{\ell}{\ell_{\mathrm{rel}}}\right)^{\alpha}.
\end{equation}

The score becomes

\begin{equation}
    s_{\mathrm{RDCC}}(D)
    = s(D)
      \left[
        1 - \chi_{s}
        \left(\frac{\ell}{\ell_{\mathrm{rel}}}\right)^{\alpha}
      \right].
\end{equation}
% ---------------------------------------------------------
\section{RDCC Diffusion Dynamics}

The RDCC reverse-time SDE is

\begin{equation}
    dD_{t}
    = - \beta(t)
      s(D_{t})
      \left[
        1 - \chi_{s}
        \left(\frac{\ell}{\ell_{\mathrm{rel}}}\right)^{\alpha}
      \right]
      dt
      + \sqrt{\beta(t)} \, d\bar{B}_{t}.
\end{equation}

RDCC induces:

\begin{itemize}
    \item reduced score magnitude at small scales,
    \item enhanced score magnitude at large scales,
    \item modified diffusion trajectories,
    \item altered generative flows,
    \item scale-dependent deformation of reverse-time dynamics.
\end{itemize}
% ---------------------------------------------------------
\section{Conclusion}

Score-based diffusion models provide a modern and flexible framework for modeling 
the distribution of persistence diagrams. In the Relaxation-Driven Cyclic 
Cosmology, the relaxation mechanism modifies the score field and the associated 
diffusion dynamics in a scale-dependent manner, producing distinctive signatures 
in weak-lensing topology. The present Companion establishes the foundation for 
future studies of generative modeling, neural SDEs, and diffusion-based inference 
in cosmological topology.

\bibliographystyle{apsrev4-2}
\nocite{*}
\bibliography{references}

\end{document}
